\section{The Learning City Point}\label{sec:learning-point}

The 2016 FTG submission introduces the concept of a ``Learning City Point''
as a one-line summary of the project's value proposition:

\begin{quote}
We'll create opportunities for participants to learn music modules in a
peeragogical style.
\end{quote}

The phrase is placed under the heading ``What is the proposed solution?'' and
is signed as a statement of intent.  No further elaboration of the Learning
City Point concept appears in the 2016 document.  The term is not defined,
nor is it linked to an existing body of literature in that document.

The present work interprets the Learning City Point as a minimal civic
infrastructure node: a publicly accessible station that lowers the barrier
to music learning by combining physical access (a keyboard, screen, and
audio system) with digital access (OER, tutoring marketplace, peer
community).  This interpretation is marked as a contemporary development,
not as a restatement of the 2013--2016 source.

\paragraph*{Relationship to the broader peeragogy framework.}
The source cites \emph{The Peeragogy V3.3} and references the peeragogy.org
community.  The Learning City Point can therefore be read as an applied
instantiation of peeragogical principles---co-learning, co-generation,
recursive practice---in a specific physical setting.  The present document
does not yet develop this into a theory; that work is reserved for a later
stage of the editorial process.
